\documentclass[11pt,a4paper]{article}
\usepackage[T1]{fontenc}
\usepackage[utf8]{inputenc}
\usepackage{lmodern}
\usepackage[margin=27mm,headheight=14pt]{geometry}
\usepackage{amsmath,amssymb,amsthm,mathtools}
\usepackage{microtype,booktabs,longtable,array}
\usepackage[dvipsnames]{xcolor}
\usepackage{listings}
\usepackage{fancyhdr}
\usepackage[hidelinks]{hyperref}
\usepackage{enumitem}
\usepackage{graphicx}
\setlength{\parskip}{4pt}
\setlength{\parindent}{1.2em}
\setlist{itemsep=2pt,topsep=4pt}
\newcommand{\R}{\mathbb R}
\newcommand{\N}{\mathbb N}
\newcommand{\OO}{\mathcal O}
\newcommand{\dd}{\partial_L}
\newcommand{\abs}[1]{\left|#1\right|}
\newcommand{\val}{\operatorname{val}}
\newcommand{\supp}{\operatorname{supp}}
\newcommand{\rising}[2]{\left(#1\right)^{\overline{#2}}}
\newcommand{\falling}[2]{\left(#1\right)^{\underline{#2}}}
\newcommand{\pkg}{\texttt{RealLogPowerInverse}}
\newtheorem{theorem}{Theorem}[section]
\newtheorem{lemma}[theorem]{Lemma}
\newtheorem{proposition}[theorem]{Proposition}
\newtheorem{corollary}[theorem]{Corollary}
\theoremstyle{definition}
\newtheorem{definition}[theorem]{Definition}
\newtheorem{remark}[theorem]{Remark}
\pagestyle{fancy}
\fancyhf{}
\fancyhead[L]{\small Real-branch log-power reversion}
\fancyhead[R]{\small Theory and implementation}
\fancyfoot[C]{\thepage}
\lstdefinestyle{wl}{basicstyle=\ttfamily\small,breaklines=true,
 columns=fullflexible,keepspaces=true,showstringspaces=false,
 frame=single,rulecolor=\color{black!20},xleftmargin=4pt,xrightmargin=4pt,
 aboveskip=8pt,belowskip=8pt,commentstyle=\color{black!55}}
\lstset{style=wl}
\hypersetup{pdftitle={Real-Branch Reversion at a Logarithmic Singularity},
 pdfauthor={OpenAI; prepared for Vladimir Reshetnikov}}
\title{\textbf{Real-Branch Reversion at a\newline Logarithmic Singularity}\\[7pt]
\large Finite log-power germs, irrational exponents,\newline
rigorous remainders, and a Wolfram Language package}
\author{Prepared for Vladimir Reshetnikov}
\date{September 7, 2026\quad---\quad Version 1.0.0}
\begin{document}
\maketitle
\begin{abstract}
We construct the positive-real inverse at the singular endpoint zero of
\(f(x)=a x+\sum_j x^{1+\alpha_j}P_j(\log x)\), where \(a>0\), the finitely many
\(\alpha_j\) are positive real numbers, and the \(P_j\) are real polynomials.
The construction answers both examples in the motivating Mathematica question:
\(x+x^2(1+\log x)\) and \(x+x^{\sqrt2}\). An auxiliary-parameter
Lagrange--B\"urmann formula is valid although ordinary Taylor reversion at zero
is not. Euler differentiation reduces every coefficient to a finite polynomial
calculation, and a locally finite additive semigroup determines exactly which
irrational powers lie below a requested cutoff. We prove local real-branch
existence, formal uniqueness, convergence for sufficiently small positive
arguments, logarithm-aware cutoff remainders, and a computable conditional
error majorant. The accompanying Wolfram Language package implements the
coefficient formula, strict input validation, exact algebraic exponent
comparison, independent sparse residual composition, and the error majorant.
It includes examples and a native regression suite. Independent Python
symbolic and numerical tests were executed; the native Wolfram suite was
prepared but could not be executed because the available evaluator failed.
\end{abstract}
\vspace{5pt}
\noindent\textbf{Principal correction to the motivating expectation.}
Writing \(L=\log y\), the first example has
\[
 g(y)=y-y^2(1+L)+y^3(2L^2+5L+3)+\OO\bigl(y^4(1+|L|)^3\bigr).
\]
Thus the two-term remainder is \(\OO(y^3\log^2 y)\), not \(\OO(y^3)\).
\newpage
\tableofcontents
\newpage

\section{The problem and the branch that must be inverted}
\subsection{What the archived question asks}
The question by Reshetnikov \cite{question} begins with the real-valued function
\begin{equation}\label{eq:original}
 f_{\mathrm{real}}(x)=
 \begin{cases}
 x+\dfrac{x^2}{2}\log(x^2e^2),&x\ne0,\\
 0,&x=0.
 \end{cases}
\end{equation}
For positive \(x\), this is simply
\begin{equation}\label{eq:logexample}
 f_1(x)=x+x^2(1+\log x).
\end{equation}
An attempted combination of \texttt{InverseFunction} and \texttt{Series}
produces expressions associated with another inverse branch, rather than the
branch tending to zero. The update to the question asks also about
\begin{equation}\label{eq:irrationalexample}
 f_2(x)=x+x^{\sqrt2},\qquad x>0.
\end{equation}
In both cases the requested answer is a germ: a function \(g(y)\) for small
positive \(y\), satisfying \(f(g(y))=y\) and \(g(y)/y\to1\).
It is not an unrestricted global complex inverse.

The two obstructions are different but have a common solution. The logarithm
prevents analyticity at zero. The irrational power prevents a usual power
series or a Puiseux series with a fixed rational denominator. Neither
obstruction prevents a convergent expansion in suitably chosen small
parameters, nor a finite calculation to any prescribed exponent cutoff.

The answer by Michael E2 in the archived discussion introduces an auxiliary
perturbation parameter, and the answer by Somos illustrates symbolic
logarithmic reversion. We use this auxiliary-parameter idea as a starting
point, but supply a branch theorem, an all-orders coefficient formula,
an exact support algorithm, and remainders. The implementation here is newly
written, not copied from those answers.

\subsection{The positive branch is not the same as all real branches}
For \eqref{eq:logexample},
\[
 f_1'(x)=1+x(3+2\log x).
\]
The function \(x(3+2\log x)\) has its minimum at \(x=e^{-5/2}\), with value
\(-2e^{-5/2}\). Consequently,
\begin{equation}\label{eq:slopeglobal}
 f_1'(x)\ge m_1:=1-2e^{-5/2}>0\qquad(x>0).
\end{equation}
Since \(f_1(0+)=0\) and \(f_1(x)\to+\infty\), there is exactly one positive
solution of \(f_1(x)=y\) for every \(y>0\). The second example is also strictly
increasing on the positive axis, since
\(f_2'(x)=1+\sqrt2\,x^{\sqrt2-1}>0\).

This does \emph{not} make \eqref{eq:original} injective on the whole real line.
For example, its negative tail tends to \(+\infty\), so nonlocal negative
solutions may coexist with the positive solution. Throughout the main
construction, the branch condition is explicitly
\begin{equation}\label{eq:branch}
 y\to0^+,
 \qquad g(y)>0,
 \qquad g(y)\sim y/a.
\end{equation}
The endpoint value \(f(0)=0\) fixes continuity; it does not supply missing
higher derivatives or select a global inverse branch for a symbolic solver.

\subsection{Why ordinary reversion is the wrong data structure}
\texttt{InverseSeries} is documented for power-series data \cite{inverseseries}.
One should not infer that an ordinary Taylor series exists merely because a
function has a first derivative at zero. In the first example, the singular
logarithm occurs in higher derivatives. In the second, the correct exponents
are \(1+n(\sqrt2-1)\), not nonnegative integers.

We therefore represent a finite answer as a list of pairs
\(\{\beta,Q_\beta(L)\}\), interpreted as

\[\sum_\beta y^\beta Q_\beta(\log y),\]
together with separate remainder metadata. A request to retain all powers \(\beta\le B\) retains each associated
logarithmic polynomial in full. In particular it does not silently turn
\(y^3\log^2 y\) into \(\OO(y^3)\).

\section{A finite log-power class and its ordered support}
\subsection{Normalization}
We study
\begin{equation}\label{eq:class}
 f(x)=a x+\sum_{j=1}^{m}x^{1+\alpha_j}P_j(\log x),
 \qquad a>0,\quad \alpha_j>0,\quad P_j\in\R[L].
\end{equation}
Zero polynomials can be deleted, and equal exponents can be merged. Put
\begin{equation}\label{eq:normalization}
 t=\frac ya,\qquad Q_j=\frac{P_j}{a},\qquad
 F(x)=\frac{f(x)}a=x+h(x),\qquad
 h(x)=\sum_j x^{1+\alpha_j}Q_j(\log x).
\end{equation}
It suffices to solve \(F(x)=t\). All coefficient formulas below refer to the
normalized \(Q_j\); the final answer is evaluated at \(t=y/a\), including
\(L=\log(y/a)\). Forgetting the factor \(a\) inside the logarithm is an error.

For \(m>0\), define
\begin{equation}\label{eq:parameters}
 \delta=\min_j\alpha_j>0,
 \qquad d=\max_j\deg Q_j,
 \qquad M(t)=1+|\log t|.
\end{equation}
As \(t\to0^+\),
\begin{equation}\label{eq:smallness}
 h(t)=\OO(t^{1+\delta}M(t)^d),\qquad
 h'(t)=\OO(t^\delta M(t)^d)=o(1).
\end{equation}
These estimates follow by differentiating each block, not by differentiating
an unspecified big-O term. If the sum is empty, the inverse is exactly \(y/a\).

\begin{proposition}[Local positive-real inverse]\label{prop:branch}
There exist \(\varepsilon,\eta>0\) such that \(F\) maps
\((0,\varepsilon)\) strictly increasingly onto \((0,\eta)\), and its inverse
\(G\) satisfies \(G(t)=t+o(t)\). The inverse of \(f\) on the branch
\eqref{eq:branch} is \(g(y)=G(y/a)\).
\end{proposition}
\begin{proof}
By \eqref{eq:smallness}, reduce \(\varepsilon\) so that
\(1/2<F'(x)<3/2\) on \((0,\varepsilon)\). Also \(F(0+)=0\).
Strict monotonicity and continuity give the inverse. From
\(F(x)/x\to1\), substituting \(x=G(t)\to0\) gives \(G(t)/t\to1\).
\end{proof}

\subsection{A locally finite, not necessarily rational, exponent grid}
Let
\begin{equation}\label{eq:semigroup}
 S=\left\{\alpha\cdot k=\sum_{j=1}^m\alpha_jk_j:
           k=(k_1,\ldots,k_m)\in\N^m\right\}.
\end{equation}
Here \(\N\) includes zero. This is an additive semigroup, rather than the full
additive group generated by the \(\alpha_j\).

\begin{lemma}[Finite support below any cutoff]\label{lem:finite}
For each \(T\ge0\), there are only finitely many multi-indices \(k\) with
\(\alpha\cdot k\le T\). In particular \(S\cap[0,T]\) is finite, and each
\(s\in S\) has only finitely many representations by such multi-indices.
\end{lemma}
\begin{proof}
Every admissible multi-index satisfies
\(0\le k_j\le\lfloor T/\alpha_j\rfloor\), and also
\(|k|\le\lfloor T/\delta\rfloor\). These are finite bounds.
\end{proof}

Irrationality does not imply accumulation at a finite positive cutoff here.
For instance \(\N+\sqrt2\N\) is locally finite, although the group
\(\mathbb Z+\sqrt2\mathbb Z\) is not. Confusing these two supports would
incorrectly suggest that a finite-order answer is impossible.

\subsection{Dominance and uniqueness of the blocks}
For any fixed \(\epsilon>0\) and integer \(r\),
\(t^\epsilon M(t)^r\to0\). Thus every logarithmic polynomial at a larger
power of \(t\) is smaller than every nonzero leading logarithmic monomial at
a smaller power. Within a fixed power, a higher logarithm degree dominates
a lower one. A nonzero polynomial has a uniquely determined leading degree.
It follows, by subtracting two expansions and inspecting their first nonzero
block, that the coefficients in a log-power expansion are unique.

The algebra used below consists of formal expressions
\(\sum_{s\in S}t^s C_s(L)\) with \(C_s\in\R[L]\). Products are well-defined
because each output weight has finitely many contributing pairs. The
valuation is the least weight of a nonzero block; the zero expression has
valuation \(+\infty\). No numerical evaluation is used to decide equality or
ordering of exponents in the package.

\section{Reversion without analyticity at zero}
\subsection{The auxiliary parameter}
For fixed positive \(t\), introduce \(\lambda\) and solve
\begin{equation}\label{eq:homotopy}
 X+\lambda h(X)=t,\qquad X(t,0)=t.
\end{equation}
The functions \(X^{1+\alpha_j}\) and \(\log X\) are holomorphic in a small
complex disk centered at \(t>0\), using the logarithm inherited from the
positive axis. Classical Lagrange--B\"urmann inversion applies there
\cite[\S1.10(vii)]{dlmf}; it is applied at \(X=t\), not at the singular point
\(X=0\). It gives
\begin{equation}\label{eq:lagrange}
 X(t,\lambda)=t+
 \sum_{n\ge1}\frac{(-\lambda)^n}{n!}
 \frac{d^{n-1}}{dt^{n-1}}h(t)^n.
\end{equation}
The derivative in this formula differentiates every occurrence of \(t\),
including \(\log t\). Freezing the logarithm while reverting an ordinary
polynomial would lose these derivative terms.

\subsection{A convergence theorem and an explicit majorant}
Fix \(0<r<1\), and let \(c_r=-\log(1-r)>0\). Write
\(Q_j(L)=\sum_{k=0}^{d_j}q_{jk}L^k\), and define
\begin{equation}\label{eq:majorant}
 \mathcal Q_r(t)=\frac1r\sum_{j=1}^{m}
 t^{\alpha_j}(1+r)^{1+\alpha_j}
 \sum_{k=0}^{d_j}|q_{jk}|\bigl(|\log t|+c_r\bigr)^k.
\end{equation}
This is a sufficient-condition majorant, not generally a sharp convergence
boundary. Since every \(\alpha_j>0\), it tends to zero as \(t\to0^+\).

\begin{theorem}[Analytic validity of auxiliary-parameter reversion]
\label{thm:convergence}
For all sufficiently small positive \(t\), the series \eqref{eq:lagrange}
converges absolutely at \(\lambda=1\), and its sum is the branch \(G(t)\)
of Proposition~\ref{prop:branch}. More generally, at any positive \(t\)
with \(\mathcal Q_r(t)<1\), it sums to the unique root in
\(|X-t|<rt\), selected by continuation from \(\lambda=0\).
This root is real and positive. For the locally defined branch, and for
\[
 G_N(t)=t+\sum_{n=1}^{N}\frac{(-1)^n}{n!}D_t^{n-1}h(t)^n,
\]
one has the explicit bound
\begin{equation}\label{eq:homotopy-bound}
 |G(t)-G_N(t)|\le
 \frac{t r}{N+1}\,
 \frac{\mathcal Q_r(t)^{N+1}}{1-\mathcal Q_r(t)}.
\end{equation}
The same bound holds at every point with \(\mathcal Q_r(t)<1\) for the
just-described disk-selected root. Without a separate global monotonicity or
continuation argument, that statement is not a claim about a globally
continued origin branch outside its local domain.
The expansion of each degree \(n\) into its finitely many multi-index
contributions is absolutely summable over all multi-indices as well.
\end{theorem}
\begin{proof}
Set \(X=t(1+u)\), with \(|u|\le r\). The branch of \(\log(1+u)\) obtained
from its Taylor series satisfies
\[
 |\log(1+u)|\le\sum_{k\ge1}\frac{r^k}{k}=c_r.
\]
Since all \(\alpha_j\) are real,
\(|(1+u)^{1+\alpha_j}|=|1+u|^{1+\alpha_j}\le(1+r)^{1+\alpha_j}\).
Consequently, on the circle \(|X-t|=rt\),
\[
 \sum_j\sup|X^{1+\alpha_j}Q_j(\log X)|
 \le rt\,\mathcal Q_r(t).
\]
For \(|\lambda|\mathcal Q_r(t)<1\), Rouch\'e's theorem compares
\(X-t+\lambda h(X)\) with \(X-t\). There is exactly one zero inside this
disk, counting multiplicities, hence it is simple. The implicit-function
theorem continues that zero holomorphically in \(\lambda\) on the disk
\(|\lambda|<1/\mathcal Q_r(t)\). The Taylor coefficients at zero are
\eqref{eq:lagrange}. For real \(\lambda=1\), complex conjugation and uniqueness
make the zero real, and \(X/t\in(1-r,1+r)\) makes it positive. For sufficiently
small \(t\), this is precisely the local branch already identified.

Let \(M_j\) denote the individual suprema just used, and put
\(M=\sum_jM_j\). Cauchy's derivative estimate gives
\[
 \left|\frac1{n!}D_t^{n-1}h(t)^n\right|
 \le\frac{(n-1)!}{n!(rt)^{n-1}}M^n
 \le\frac{rt}{n}\mathcal Q_r(t)^n.
\]
Summing for \(n>N\), and replacing \(1/n\) by \(1/(N+1)\), proves
\eqref{eq:homotopy-bound}.
For absolute summability before merging multi-indices, apply Cauchy's
estimate to \(\prod_jh_j^{k_j}\) and sum over \(|k|=n\). The multinomial
identity gives
\[
 \frac{(n-1)!}{(rt)^{n-1}}
 \sum_{|k|=n}\frac{\prod_jM_j^{k_j}}{\prod_j k_j!}
 =\frac{M^n}{n(rt)^{n-1}}.
\]
The same summable geometric majorant results. Thus regrouping by exponent is
legitimate, even when different multi-indices have the same weight.
\end{proof}

\begin{remark}[What ``convergent'' means here]
The series is not an ordinary Taylor series in \(t\). It is a convergent
log-power series evaluated for sufficiently small positive \(t\), and it is
also an asymptotic expansion in the ordered log-power scale. Neither claim
asserts analyticity at zero. Convergence of the auxiliary series does not
imply that a default branch selected by \texttt{InverseFunction} is the
branch being represented.
\end{remark}

\section{The finite coefficient formula}
\subsection{Euler differentiation of a block}
For every real \(\beta\) and polynomial \(P\),
\begin{equation}\label{eq:euler}
 D_t\bigl(t^\beta P(L)\bigr)
 =t^{\beta-1}(\beta+\dd)P(L),\qquad L=\log t.
\end{equation}
Repeated differentiation gives
\begin{equation}\label{eq:repeated-euler}
 D_t^r\bigl(t^\beta P(L)\bigr)
 =t^{\beta-r}\prod_{j=0}^{r-1}(\beta-j+\dd)P(L).
\end{equation}
All factors commute, since they are constant-coefficient polynomials in
\(\dd\). Their application never increases the logarithm degree.

For \(k\in\N^m\setminus\{0\}\), write
\[
 n=|k|,\qquad s=\alpha\cdot k,\qquad
 k!=\prod_j k_j!,\qquad Q^k(L)=\prod_j Q_j(L)^{k_j}.
\]
Expanding \(h^n\) by the multinomial theorem and using
\eqref{eq:repeated-euler} yields the central formula.

\begin{theorem}[Multi-index inversion formula]\label{thm:coefficients}
For the branch in \eqref{eq:branch},
\begin{equation}\label{eq:master}
 \boxed{\displaystyle
 g(y)=t+\sum_{k\in\N^m\setminus\{0\}}
 t^{1+\alpha\cdot k}\,
 \frac{(-1)^{|k|}}{k!}
 \rising{\dd+2+\alpha\cdot k}{|k|-1}
 \prod_jQ_j(L)^{k_j},\quad
 t=\frac ya,\ L=\log t.}
\end{equation}
An empty operator product equals the identity. This formula is formally
well-defined coefficientwise; it is absolutely convergent for sufficiently
small positive \(y\). Its power support is contained in \(1+S\), and the
contribution from \(k\) has logarithm degree at most
\(\sum_jk_j\deg Q_j\).
\end{theorem}
\begin{proof}
The term of \(h(t)^n\) indexed by \(k\) is
\(n!\,t^{n+s}Q^k(L)/k!\). Differentiating \(n-1\) times lowers the power to
\(1+s\), and the operator factors are
\(\dd+n+s,\ldots,\dd+2+s\). They are exactly the rising product in
\eqref{eq:master}. The factor \(n!\) cancels that in \eqref{eq:lagrange}.
Local finiteness follows from Lemma~\ref{lem:finite}, and analytic validity
from Theorem~\ref{thm:convergence}.
\end{proof}

\subsection{A completely scalar coefficient formula}
Define numbers \(r_{n,s,j}\) by
\[
 \rising{z+2+s}{n-1}=\sum_{j=0}^{n-1}r_{n,s,j}z^j.
\]
The block contributed by \(k\) is therefore
\begin{equation}\label{eq:scalar}
 C_k(L)=\frac{(-1)^n}{k!}
 \sum_{j=0}^{\min(n-1,\deg Q^k)}r_{n,s,j}\,(Q^k)^{(j)}(L).
\end{equation}
This provides a direct formula in finite products and polynomial
coefficients, without symbolic differentiation of irrational powers. The
numbers \(r_{n,s,j}\) can equivalently be expressed through coefficients of
rising factorials; that additional notation is not needed in the code.

For a single block \(h=t^{1+\alpha}cL^q\), it specializes to
\[
 C_n(L)=\frac{(-c)^n}{n!}
 \sum_{j=0}^{\min(n-1,nq)}r_{n,n\alpha,j}
 \falling{nq}{j}L^{nq-j}.
\]
For multiple generators, coincident weights must be \emph{added}, not
selected arbitrarily. These coincidences are the only ``resonances'' needed
in this finite positive-exponent problem. There is no small-denominator
inversion in \eqref{eq:master}.

\section{Cutoffs, exact gaps, and rigorous remainders}
\subsection{What the cutoff requests}
For an absolute power cutoff \(B\ge1\), set \(T=B-1\), and let \(g_B\) be
\eqref{eq:master} restricted to \(\alpha\cdot k\le T\), together with the
leading term \(t\). The least omitted semigroup point is
\begin{equation}\label{eq:nextgap}
 s_+(T)=\min\{s\in S:s>T\}.
\end{equation}
It exists because \(S\) contains \(\delta\N\), and it is computable from a
finite enumeration: it lies in \((T,T+\delta]\).

\begin{theorem}[Logarithm-aware exponent-cutoff remainder]
\label{thm:remainder}
With \(s_+=s_+(T)\),
\begin{equation}\label{eq:remainder}
 g(y)-g_B(y)=
 \OO\left(t^{1+s_+}M(t)^{d\lfloor s_+/\delta\rfloor}\right)
 \qquad(t=y/a\to0^+).
\end{equation}
The exponent and logarithm degree in this bound are conservative if omitted
coefficient blocks cancel. No nonzero-coefficient assumption is needed.
\end{theorem}
\begin{proof}
Choose an integer \(N\) with \((N+1)\delta>s_+\).
The tail of the auxiliary-parameter series after degree \(N\), by
Theorem~\ref{thm:convergence} and \(\mathcal Q_r(t)=\OO(t^\delta M^d)\), is
\[
 \OO\bigl(t^{1+(N+1)\delta}M^{d(N+1)}\bigr)
 =o(t^{1+s_+}).
\]
Every retained multi-index has \(|k|\delta\le T<s_+\), hence \(|k|\le N\).
Among the finitely many degree-\(\le N\) contributions not retained, weights
strictly larger than \(s_+\) are \(o(t^{1+s_+})\), even after multiplication
by their fixed logarithmic polynomials. At weight exactly \(s_+\),
\(|k|\le\lfloor s_+/\delta\rfloor\), so the logarithm degree is at most
\(d\lfloor s_+/\delta\rfloor\). Adding these finite contributions proves
\eqref{eq:remainder}.
\end{proof}

\subsection{An explicit numerical-point bound for an exponent cutoff}
The estimate above gives an asymptotic scale, not a constant or an endpoint
threshold. A different, usually looser, bound supplies an explicit constant.
Put \(\alpha_{\max}=\max_j\alpha_j\) and
\begin{equation}\label{eq:safedegree}
 N_0=\left\lfloor\frac{B-1}{\alpha_{\max}}\right\rfloor.
\end{equation}
Every multi-index of total degree \(n\le N_0\) is retained, since
\(\alpha\cdot k\le n\alpha_{\max}\le B-1\).
The omitted indices form a subset of the degrees \(n>N_0\). Absolute
multi-index summability in Theorem~\ref{thm:convergence} therefore gives
\begin{equation}\label{eq:cutoff-bound}
 \boxed{\displaystyle
 |g(y)-g_B(y)|\le
 \frac{t r\,\mathcal Q_r(t)^{N_0+1}}
 {(N_0+1)(1-\mathcal Q_r(t))}}
 \quad\text{provided }\mathcal Q_r(t)<1.
\end{equation}
This is the inequality implemented by \texttt{InverseErrorBound}. At an
arbitrary evaluation point, the selected root is specified by
\(t(1-r)<X<t(1+r)\); near zero it is the requested germ. If the input is
not globally monotone, using the bound for a different global branch requires
additional justification. The two original examples are globally monotone
on the positive axis, so this distinction causes no ambiguity there. It can be
weak when the generators have very different sizes. It is nevertheless a
valid bound for the requested exponent cutoff, not just for a homotopy
truncation. A floating-point evaluation of its right-hand side is not an
outward-rounded interval certificate; the package reports the condition
explicitly and makes no such numerical certification claim.

\subsection{Finite determination and remainder perturbations}
The coefficient at weight \(s\) depends only on input blocks of weights
\(\le s\). Hence a finite input jet determines a finite inverse jet.
An analytic version requires a condition on the actual remainder, not only
formal notation.

\begin{proposition}[Stability under a controlled input remainder]
\label{prop:stability}
Suppose \(F_0=x+h(x)\) is as above, and
\(F(x)=F_0(x)+R(x)\), with
\[
 R(x)=\OO(x^{1+\rho}M(x)^q),\qquad R'(x)=o(1),\qquad\rho>0.
\]
Then both local positive inverses exist, and
\[
 F^{-1}(t)-F_0^{-1}(t)
 =\OO(t^{1+\rho}M(t)^q).
\]
\end{proposition}
\begin{proof}
Both derivatives tend to one, so \(F'\ge1/2\) after reducing the interval.
Let \(x_0=F_0^{-1}(t)\sim t\). The residual of \(x_0\) in the perturbed
equation is \(R(x_0)\). The mean-value theorem gives
\(|F^{-1}(t)-x_0|\le2|R(x_0)|\), and \(x_0\sim t\) transports the bound.
\end{proof}
This proposition also explains why a finite log-power expansion cannot
identify exponentially small additions such as \(e^{-1/x}\). Such additions
may leave every algebraic-logarithmic coefficient unchanged while changing
the actual inverse by a flat term. They are not part of the finite exact input
class accepted by this version of the package.

\section{Formal composition and a second construction}
\subsection{The normalized correction equation}
Write \(G(t)=t(1+u)\). Its equation is
\begin{equation}\label{eq:correction}
 u+\sum_j t^{\alpha_j}(1+u)^{1+\alpha_j}
 Q_j\bigl(L+\log(1+u)\bigr)=0.
\end{equation}
Since \(u\) has strictly positive valuation, both
\begin{equation}\label{eq:jets}
 (1+u)^r=\sum_{n\ge0}\binom rn u^n,
 \qquad
 \log(1+u)=\sum_{n\ge1}\frac{(-1)^{n+1}}n u^n
\end{equation}
are well-defined formal series on the same semigroup. Only finitely many
powers contribute to a fixed cutoff. Substitution into a polynomial
\(Q_j\) is likewise finite after truncation.

\begin{proposition}[Formal uniqueness and coefficient stabilization]
\label{prop:formal}
The map
\[
 \mathcal T(u)=-\sum_jt^{\alpha_j}(1+u)^{1+\alpha_j}
 Q_j\bigl(L+\log(1+u)\bigr)
\]
has a unique positive-valuation fixed point. The iterates \(u_0=0\),
\(u_{n+1}=\mathcal T(u_n)\), stabilize coefficientwise; after \(n\) iterations,
\(\val(u-u_n)\ge(n+1)\delta\).
\end{proposition}
\begin{proof}
For positive-valuation \(u,v\), subtract their formal binomial and logarithm
series. Each difference is divisible by \(u-v\), with quotient of
nonnegative valuation. The explicit factor \(t^{\alpha_j}\) consequently
implies
\[
 \val(\mathcal T(u)-\mathcal T(v))\ge\val(u-v)+\delta.
\]
Starting from zero gives \(\val(u_1-u_0)\ge\delta\), and successive
differences gain at least \(\delta\). By local finiteness, every coefficient
stabilizes. The limit is a fixed point. Two different fixed points would have
a finite first difference weight \(s\), but the displayed inequality would
force that same weight to be at least \(s+\delta\), a contradiction.
\end{proof}

Thus the Lagrange formula and fixed-point construction must give the same
formal answer. This supplies an independent way to check coefficients and a
conceptually simple alternative implementation.

\subsection{Residual validation is not numerical root finding}
For a finite correction \(u_B\), the residual checker evaluates the left side
of \eqref{eq:correction} by sparse multiplication and the finite jets
\eqref{eq:jets}. It does not call the Euler-operator formula that generated
the answer. If the result vanishes through weight \(B-1\), then
\(f(g_B(y))-y\) vanishes through absolute power \(B\).

This statement is exact finite algebra. It does not automatically certify a
floating-point root approximation. To obtain an actual error estimate from a
numerical residual, one also needs a derivative lower bound and knowledge
that the approximation and desired root lie in the chosen branch interval.
If \(f'\ge m>0\) there, then
\begin{equation}\label{eq:residual-bound}
 |g_B(y)-g(y)|\le\frac{|f(g_B(y))-y|}{m}.
\end{equation}
A bracket \(f(\ell)\le y\le f(r)\) on a strictly increasing branch gives
\(g(y)\in[\ell,r]\). For rigorous numerical use these inequalities must be
checked exactly or with outward-rounded intervals.

For \(f_1\), the global bound \eqref{eq:slopeglobal} supplies \(m=m_1\) for
any positive approximation; for \(f_2\), one may use \(m=1\).
These are stronger practical tools than a small raw residual alone.

\section{The logarithmic example in full detail}
\subsection{All-order polynomial blocks}
For \(f_1(x)=x+x^2(1+\log x)\), there is one generator \(\alpha=1\), and
\begin{equation}\label{eq:logseries}
 g_1(y)=y+\sum_{n\ge1}y^{n+1}C_n(L),\qquad L=\log y,
\end{equation}
where
\begin{equation}\label{eq:logcoeff}
 C_n(L)=\frac{(-1)^n}{n!}
 \rising{\dd+n+2}{n-1}(1+L)^n.
\end{equation}
The first four polynomials are
\begin{align}
 C_1(L)&=-L-1,\nonumber\\
 C_2(L)&=2L^2+5L+3,\nonumber\\
 C_3(L)&=-5L^3-\frac{41}{2}L^2-27L-\frac{23}{2},\label{eq:logpolys}\\
 C_4(L)&=14L^4+\frac{241}{3}L^3+\frac{335}{2}L^2+151L+\frac{299}{6}.
 \nonumber
\end{align}
The next block, useful as an independent test fixture, is
\begin{equation}\label{eq:c5}
 C_5(L)=-42L^5-\frac{3727}{12}L^4-\frac{5365}{6}L^3
        -1256L^2-\frac{5177}{6}L-\frac{2789}{12}.
\end{equation}
For example, the \(n=2\) Lagrange term is
\[
 \frac12D_y\bigl(y^4(1+L)^2\bigr)
 =y^3\bigl(2(1+L)^2+(1+L)\bigr),
\]
which explains the \(5L+3\) terms that would be lost by incorrectly treating
\(L\) as a constant during differentiation.

\subsection{Catalan leading logarithms and the exact remainder scale}
Only the zero-derivative part of the operator contributes to degree \(n\)
in \(C_n\), so
\begin{equation}\label{eq:catalan}
 [L^n]C_n(L)=\frac{(-1)^n}{n!}\rising{n+2}{n-1}
 =(-1)^n\frac1{n+1}\binom{2n}{n}.
\end{equation}
These are signed Catalan numbers, and they are nonzero. It follows that
\begin{equation}\label{eq:logremainder}
 g_1(y)-\left(y+\sum_{n=1}^{N}y^{n+1}C_n(\log y)\right)
 =\OO\bigl(y^{N+2}M(y)^{N+1}\bigr),
\end{equation}
with that logarithmic degree actually present in the first omitted block.
In particular,
\[
 \frac{g_1(y)-y+y^2(1+\log y)}{y^3}
 =2\log^2y+5\log y+3+o(1),
\]
where the \(o(1)\) follows from the next remainder divided by \(y^3\).
The ratio is unbounded, proving that \(\OO(y^3)\) is not the correct
remainder for the two displayed leading terms.

\subsection{An analytic two-parameter interpretation}
The correction equation becomes
\[
 u+y(1+u)^2\bigl(1+L+\log(1+u)\bigr)=0.
\]
Introduce \(v=y(1+L)\) and \(w=y\). This is
\begin{equation}\label{eq:two-parameter}
 u+(1+u)^2\bigl(v+w\log(1+u)\bigr)=0.
\end{equation}
The analytic implicit-function theorem at \((u,v,w)=(0,0,0)\) gives a
convergent two-variable power series \(u=U(v,w)\). Since both
\(y(1+\log y)\) and \(y\) tend to zero, evaluating this analytic function
along that curve produces \eqref{eq:logseries}. The logarithmic singularity
of the curve is compatible with analyticity of the lifted function.

A sharper one-example majorant than \eqref{eq:majorant} is obtained by using
\(|1+L|\) directly. For \(r=1/2\), set
\begin{equation}\label{eq:logq}
 q(y)=\frac92 y\bigl(|1+\log y|+\log2\bigr).
\end{equation}
If \(q(y)<1\), Cauchy's estimate gives
\begin{equation}\label{eq:logbound}
 |g_1(y)-g_{1,N}(y)|\le
 \frac{y}{2(N+1)}\frac{q(y)^{N+1}}{1-q(y)}.
\end{equation}
This is a sufficient, explicitly checkable domain, not a claim that the
series ceases to converge whenever the inequality fails. The generic package
majorant is slightly looser because it sums absolute polynomial coefficients.

\subsection{The negative-side germ of the original real function}
For \(x=-u\), \(u>0\), and \(y=-v\), the equation
\(f_{\mathrm{real}}(x)=y\) is
\[
 v=u-u^2(1+\log u).
\]
It is another member of our class, with polynomial \(-(1+L)\). Substituting
back shows that the local two-sided inverse can be written with the same
integer powers and with \(\log|y|\):
\[
 g_{\mathrm{real}}(y)=y+\sum_{n\ge1}y^{n+1}C_n(\log|y|),\qquad y\ne0
\]
for sufficiently small \(|y|\). To compute the negative side using the
positive-only package, invert \(u-u^2(1+\log u)\) for \(v>0\) and then set
\(y=-v\), \(g(y)=-u(v)\). The package never globally rewrites
\(\log(x^2)\) as \(2\log x\) without the positive-real assumption.

\section{The irrational-power example and its convergence radius}
\subsection{A one-generator generalized power series}
Consider the entire family
\(f_p(x)=x+x^p\), with \(p>1\), and put \(\alpha=p-1\). The inverse is
\begin{equation}\label{eq:pseries}
 g_p(y)=\sum_{n\ge0}c_n(p)y^{1+n\alpha},\qquad c_0(p)=1,
\end{equation}
where the main theorem gives, for \(n\ge1\),
\begin{align}
 c_n(p)
 &=\frac{(-1)^n}{n!}\prod_{j=0}^{n-2}(np-j)\label{eq:pcoefficient}\\
 &=\frac{(-1)^n\Gamma(np+1)}
 {\Gamma(n+1)\Gamma(n(p-1)+2)}
 =\frac{(-1)^n}{1+n(p-1)}\binom{np}{n}.\nonumber
\end{align}
The empty product for \(n=1\) is one. The finite product, rather than gamma
quotients, is used for symbolic coefficient generation in the package.

The first terms for general \(p\) are
\begin{align}\label{eq:pgeneral}
 g_p(y)={}&y-y^p+p\,y^{2p-1}
 -\frac{p(3p-1)}2y^{3p-2}\\
 &+\frac{p(4p-1)(4p-2)}6y^{4p-3}
 +\OO(y^{5p-4}).\nonumber
\end{align}
For the requested value \(p=\sqrt2\), this becomes
\begin{align}\label{eq:sqrt2}
 g_2(y)={}&y-y^{\sqrt2}+\sqrt2\,y^{2\sqrt2-1}
 -\left(3-\frac{\sqrt2}{2}\right)y^{3\sqrt2-2}\\
 &+\left(\frac{17\sqrt2}{3}-4\right)y^{4\sqrt2-3}
 +\OO(y^{5\sqrt2-4}).\nonumber
\end{align}
There are finitely many terms below every exponent cutoff. To retain terms
through the \(n=N\) contribution, pass the exact cutoff
\(B=1+N(\sqrt2-1)\), not a decimal approximation to it.

\subsection{Why this is not a Puiseux series}
If \(p\) is rational, a suitable root of \(y\) makes all exponents integral.
If \(p\) is irrational, no fixed rational denominator does so. Nevertheless,
\(z=y^{p-1}\) makes the correction an ordinary analytic series:
\[
 g_p(y)=yW(z),\qquad W(z)+zW(z)^p=1,\qquad W(0)=1.
\]
The implicit derivative with respect to \(W\) is one at \((W,z)=(1,0)\),
which immediately proves a nonzero radius of convergence in \(z\).

\subsection{Exact radius from the coefficient formula}
Stirling's formula applied to \eqref{eq:pcoefficient} gives
\begin{equation}\label{eq:pasympt}
 |c_n(p)|\sim
 \frac{\sqrt p}{\sqrt{2\pi}(p-1)^{3/2}}\,
 n^{-3/2}\rho^{-n},
 \qquad
 \rho=\frac{(p-1)^{p-1}}{p^p}.
\end{equation}
Indeed, the exponential factors cancel between the three gamma functions;
the remaining exponential base is \(p^p/(p-1)^{p-1}\), and the power of \(n\)
is \(-3/2\). Thus the radius of \(W(z)\) is exactly \(\rho\). On the positive
axis, \eqref{eq:pseries} certainly converges for
\begin{equation}\label{eq:pthreshold}
 0<y<\rho^{1/(p-1)}.
\end{equation}
The coefficient estimate also gives absolute convergence at \(|z|=\rho\),
although analyticity need not extend through the negative boundary point.
For \(p=\sqrt2\), numerical evaluation gives
\[
 \rho\approx0.4251956066518408138025026,\qquad
 \rho^{1/(p-1)}\approx0.1268627051233096813979122.
\]
The series boundary is not a boundary of existence of the positive inverse:
that inverse exists for every \(y>0\). It is a boundary for this expansion
about the chosen auxiliary coordinate \(z=0\).

\section{Mixed generators and coincident exponents}
\subsection{The first mixed term}
For
\[
 F(x)=x+x^{1+\alpha}P(\log x)+x^{1+\beta}Q(\log x),
\]
the contributions with multi-indices \((1,0),(0,1),(1,1)\) are
\begin{align*}
 -t^{1+\alpha}P(L),\qquad
 -t^{1+\beta}Q(L),\qquad
 t^{1+\alpha+\beta}(\dd+2+\alpha+\beta)(P(L)Q(L)).
\end{align*}
The two pure quadratic contributions are
\[
 \frac12t^{1+2\alpha}(\dd+2+2\alpha)P(L)^2,
 \qquad
 \frac12t^{1+2\beta}(\dd+2+2\beta)Q(L)^2.
\]
These terms must be sorted by their actual exponents, not by total
multi-index degree. For example, a contribution of degree three in a small
generator can precede a contribution of degree one in a large generator.

\subsection{A resonant worked example}
Take \(\alpha=1/2\), \(\beta=1\),
\(P(L)=1+L\), and \(Q(L)=2-L\). The exponent \(1+\beta=2\) coincides with
\(1+2\alpha=2\); both contributions are required. The result is
\begin{align}\label{eq:resonant}
 G(t)={}&t-(1+L)t^{3/2}
 +\left(\frac32L^2+5L+\frac12\right)t^2\\
 &+\left(-\frac{21}{8}L^3-\frac{123}{8}L^2
         -\frac{123}{8}L+\frac38\right)t^{5/2}
 +\OO(t^3M(t)^4).\nonumber
\end{align}
The remainder shown uses the conservative general theorem. The package
merges equal input generators before enumeration and equal output exponents
afterward. It also deletes blocks whose combined polynomial is zero.

Even an ordinary polynomial exposes the cancellation issue. For
\(F=x+x^2+2x^3\), the inverse has \(t-t^2+0t^3+\cdots\), whereas
\(F=x+x^2+3x^3\) has \(t-t^2-t^3+\cdots\). A data structure that overwrites a
coincident exponent rather than adding its coefficients fails these tests.

\section{The Wolfram Language implementation}
\subsection{Supported contract and public entry points}
The mathematical theory permits arbitrary fixed positive real generators.
The software deliberately restricts generators and the cutoff to exact real
algebraic numbers. This makes ordering, equality, and bounded lattice
enumeration effectively decidable with \texttt{RootReduce} and exact
symbolic comparisons \cite{rootreduce,algebraics}. Polynomial coefficients
and the positive linear coefficient may be exact, provably real numeric
constants, such as rationals, radicals, or \(\pi\); unresolved symbolic
parameters and machine-precision constants are not accepted.

The main public interfaces are:
\begin{lstlisting}
RealInverseExpansion[f, {x,y}, B]
LogPowerInverseExpansion[{{alpha1,P1}, ...}, {y,L}, B]
InverseResidual[data, B]
InverseErrorBound[data, y0]
PurePowerInverseCoefficient[p, n]
LogQuadraticInversePolynomial[n, L]
\end{lstlisting}
The structured interface accepts the blocks of the original, unnormalized
\(f\): \(x^{1+\alpha_j}P_j(\log x)\). Its option
\texttt{"LinearCoefficient" -> a} performs the division by \(a\) internally.
The supplied logarithm symbol \texttt{L} must differ from the output symbol.
The expression front end normalizes under \(x>0\) and then parses finite
monomials. It never uses an unrestricted \texttt{PowerExpand}.

Clear the symbols used as variables before calling the package. As with
ordinary Wolfram Language algebra, pre-existing values of \texttt{x},
\texttt{y}, or the logarithm symbol can alter an expression before it reaches
the function. The documented interfaces accept expressions, not held
unevaluated programs.

\subsection{Loading and the two original examples}
From the extracted archive root:
\begin{lstlisting}
Get["Kernel/RealLogPowerInverse.wl"];
Clear[x,y,L];

r1 = RealInverseExpansion[x+x^2 (1+Log[x]), {x,y}, 5];
r1["Expression"]
r1["Remainder"]
InverseResidual[r1]["Vanishes"]

r2 = RealInverseExpansion[x+x^Sqrt[2],
                         {x,y}, 1+6 (Sqrt[2]-1)];
r2["Expression"]
InverseResidual[r2]["Vanishes"]
\end{lstlisting}
The residual outputs are expected to be \texttt{True}; these exact calls are
included in the native regression suite. They are not represented here as
observed native-kernel output, because that suite was not executed in this
environment.

The structured equivalent of the first call is
\begin{lstlisting}
r1 = LogPowerInverseExpansion[{{1,1+L}}, {y,L}, 5];
\end{lstlisting}
For the original \texttt{Piecewise} definition, one can first simplify the
positive restriction:
\begin{lstlisting}
f = Piecewise[{{x+x^2 Log[x^2 E^2]/2, x!=0}}, 0];
fPositive = FullSimplify[f, Assumptions -> x>0];
RealInverseExpansion[fPositive, {x,y}, 3]
\end{lstlisting}
The explicit structured interface remains available when the expression
front end cannot normalize a semantically equivalent input form.

\subsection{Returned data and remainder semantics}
Each expansion is an \texttt{Association}. The most important fields are
\begin{center}
\begin{tabular}{@{}p{0.28\linewidth}p{0.65\linewidth}@{}}
\toprule
Field & Meaning\\\midrule
\texttt{"Expression"} & The finite expression in \(y/a\) and \(\log(y/a)\).\\
\texttt{"Terms"} & Pairs \(\{\beta,C(L)\}\), with absolute powers \(\beta\).\\
\texttt{"CorrectionBlocks"} & Pairs \(\{s,C(L)\}\) representing \(G(t)/t-1\).\\
\texttt{"InputBlocks"} & Normalized pairs \(\{\alpha_j,P_j(L)/a\}\).\\
\texttt{"Remainder"} & Power, logarithm degree, scale, and the asymptotic interpretation.\\
\texttt{"MultiIndexCount"} & Number enumerated, including the zero index and the finite next-gap lookahead.\\
\bottomrule
\end{tabular}
\end{center}
The remainder scale is an ordinary expression accompanied by a statement of
its big-O meaning. It is not a \texttt{SeriesData} object with an invented
integer order. Its implicit constant and threshold are not supplied.
For a computable conditional bound at an exact point, use, for example,
\begin{lstlisting}
b = InverseErrorBound[r1, 1/1000];
b["Condition"]
b["Bound"]
\end{lstlisting}
The bound is a \texttt{ConditionalExpression} with condition
\(\mathcal Q_r<1\), using \texttt{"Radius" -> 1/2} by default. Failure of
this sufficient condition is not evidence that the expansion is invalid.

\subsection{Core algorithm and termination}
The coefficient engine performs the following finite calculation. It first
validates and merges the normalized input blocks. It then enumerates
\(k\in\N^m\) with \(\alpha\cdot k\le B-1+\delta\) recursively, using the
coordinate bound
\[
 k_j\le\left\lfloor
 \frac{B-1+\delta-\sum_{i<j}\alpha_i k_i}{\alpha_j}
 \right\rfloor.
\]
The extra \(\delta\) finds the next semigroup point for the remainder.
For every retained nonzero index it constructs \(Q^k\), applies the Euler
factors, divides by \(k!\), and merges the result by exact exponent equality.
The essential polynomial loop is
\begin{lstlisting}
q = Times @@ MapThread[Power, {polynomials, k}];
Do[q = Expand[(2+s+j) q + D[q,L]], {j,0,n-2}];
coefficient = (-1)^n q / (Times @@ (Factorial /@ k));
\end{lstlisting}
A constant term in a polynomial is handled by the same code; for \(n=1\)
the loop is empty. Factorial normalization is multi-index factorial
\(\prod k_j!\), not \(|k|!\), because the multinomial coefficient has already
been accounted for.

The number of lattice points is bounded by
\(\prod_j(1+\lfloor(B-1+\delta)/\alpha_j\rfloor)\). This bound can be large;
it is not a polynomial-time promise in the number of generators. The option
\texttt{"MaxMultiIndices" -> 100000} stops excessive enumeration with a
\texttt{Failure} object. The independent residual checker has its own
\texttt{"MaxProducts" -> 1000000} guard. Exact algebraic comparison and
polynomial simplification can still be costly; these guards count algebraic
work units rather than elapsed time or total memory.

\subsection{Independent sparse composition}
The residual engine stores a sparse series as a sorted list of
\(\{s,P(L)\}\) pairs. Addition merges coefficients, multiplication adds
weights and discards terms past the cutoff, binomial and logarithm series
stop when their next power has no remaining support, and polynomial
substitution uses Horner's rule. All correction weights must be positive;
this is what ensures the loops terminate.

The result includes \texttt{"NormalizedResidual"}, representing
\((f(g_B(y))-y)/y\), and \texttt{"Residual"}, representing the unnormalized
residual after truncation. The checker verifies the polynomial data in
\texttt{"CorrectionBlocks"}; hand-editing only \texttt{"Expression"} does
not change that data and is outside the supported workflow. Treat returned
associations as immutable records, or rebuild them consistently.

\subsection{Deliberate limitations}
The implementation is not a universal transseries system. It rejects
nonpositive excess exponents, iterated logarithms, negative logarithm powers,
oscillatory coefficients, approximate or nonalgebraic exponent data, and
unspecified input remainders. It does not determine all global inverse
branches, prove arbitrary user assumptions, or formally verify itself in a
proof assistant. The two requested examples and their finite mixed-log-power
generalizations lie fully inside its supported class.

The two coefficient helper functions return formal symbolic polynomials and
are not branch validators. For example,
\texttt{PurePowerInverseCoefficient[p,n]} is useful symbolically even when
\texttt{p} is a parameter; the inverse-function theorem for that family still
requires \(p>1\). The main expansion interfaces enforce their stricter
numeric-domain contract.

\section{Verification, numerical evidence, and reproducibility}
\subsection{What was actually executed}
The independent script \texttt{validation/reference\_validation.py} was run
with SymPy and mpmath. It completed 19 grouped symbolic checks and 27
high-precision numerical cases. The symbolic checks include direct residual
composition, fixed-point cross-checks, mixed irrational and resonant support,
scaled coefficients, noninteger cutoffs, Catalan leading coefficients,
ordinary quadratic inversion, and detection of a deliberately corrupted
coefficient. Its polynomial calculations are exact, not fitted to numerical
samples.

The numerical comparisons were performed at 110 decimal digits. The reported
errors compare truncated expansions with high-precision positive roots.
They are numerical evidence, not interval-arithmetic certificates. The
logarithmic cases additionally check the numerical residual bound using
\(m_1\) from \eqref{eq:slopeglobal}.

A genuine Wolfram evaluator call failed with an unavailable server/dependency
error, and no local Wolfram kernel was present. Consequently the delivered
Wolfram regression suite was \emph{not} executed. This is an explicit release
limitation: independent mathematical validation does not establish native
Wolfram execution, parser behavior, or compatibility on every claimed target
version. The full package source and tests are supplied so those checks are
reproducible in the user's installation.

\subsection{Representative numerical results}
Here \(N\) is the number of nonleading contributions retained in
\eqref{eq:logseries} or \eqref{eq:pseries}, not an absolute exponent cutoff.
The ratio compares the absolute error with the magnitude of the first
omitted term.
\begin{center}
\begin{tabular}{@{}llrrr@{}}
\toprule
Example & \(y\) & \(N\) & Absolute error & Error/next term\\\midrule
Logarithmic & \(10^{-2}\) & 1 & \(2.42021\,10^{-5}\) & 1.08096\\
Logarithmic & \(10^{-2}\) & 6 & \(1.07915\,10^{-10}\) & 1.10472\\
Logarithmic & \(10^{-3}\) & 3 & \(1.25941\,10^{-11}\) & 1.01590\\
Logarithmic & \(10^{-6}\) & 1 & \(3.15669\,10^{-16}\) & 1.00003\\
Logarithmic & \(10^{-6}\) & 6 & \(1.79594\,10^{-38}\) & 1.00004\\\addlinespace
Irrational power & \(10^{-2}\) & 1 & \(2.51975\,10^{-4}\) & 0.80853\\
Irrational power & \(10^{-2}\) & 6 & \(3.37743\,10^{-7}\) & 0.77205\\
Irrational power & \(10^{-3}\) & 3 & \(3.88767\,10^{-8}\) & 0.90508\\
Irrational power & \(10^{-6}\) & 1 & \(1.50541\,10^{-11}\) & 0.99473\\
Irrational power & \(10^{-6}\) & 6 & \(1.09682\,10^{-22}\) & 0.99351\\
\bottomrule
\end{tabular}
\end{center}
The ratios tending to one as \(y\) decreases agree with the nonzero leading
omitted-block analysis. The irrational-power example approaches its small
parameter regime more slowly because \(y^{\sqrt2-1}\) decreases more slowly
than \(y\).

\subsection{Running the native tests and rebuilding the article}
The native suite contains 34 \texttt{VerificationTest} cases. It exercises
expression and structured interfaces, both examples, exact collisions and
cancellations, scaling, explicit error bounds, invalid input, resource limits,
and detection of corrupted correction data. Wolfram's \texttt{TestReport}
interface is documented in \cite{testreport}. From the archive root run
\begin{lstlisting}
TestReport["Tests/RealLogPowerInverse.wlt"]
\end{lstlisting}
or from a terminal with a working kernel:
\begin{lstlisting}
wolframscript -file Tests/RunTests.wls
python validation/reference_validation.py
cd article
pdflatex -interaction=nonstopmode -halt-on-error article.tex
pdflatex -interaction=nonstopmode -halt-on-error article.tex
\end{lstlisting}
The first command is intentionally a user-run native test, not a report of
successful execution here. The Python output is recorded in
\texttt{validation/results.json} and \texttt{validation/run.log}.
The bibliography is embedded in this source; no external bibliography file
or network access is required to build the PDF.

\section{Relation to the supplied research materials}\label{sec:related}
The supplied ProveIt materials organize polynomial-logarithmic scales,
well-based support, differential block operators, reversion, and remainder
transport \cite{proveit}. Their orientation and status files explicitly
distinguish classical coefficient Lagrange--B\"urmann inversion from
near-identity operator formulations, and distinguish formal identities from
analytic branch and convergence claims. Those distinctions are important
here: the passage from \eqref{eq:lagrange} to \eqref{eq:master} is proved,
not inferred from the shared name ``Lagrange inversion.''

The connection is concrete. The block derivative is
\(t^{\beta-1}(\beta+\partial_L)\); the positive semigroup supplies finite
factorization and cutoff termination; composition closes on the same support;
and a real derivative lower bound transports residuals to actual inverse
errors. The present setting is intentionally smaller than a full transseries
field. Its positive, finite generators permit elementary proofs and a compact
implementation, without importing the full machinery of well-based
transseries.

\subsection*{Scope of source inspection}
The repository reference is pinned to the following commit:
\begin{center}
\texttt{24ce8bd743eaab64a91ce90725ea00f498d319d2}.
\end{center}
The group and consolidated-volume
README files and targeted search excerpts were accessible. The oversized
consolidated \TeX{} volume could not be retrieved in full through the available
connector. Accordingly, this article does not claim a line-by-line review of
that volume or verification of its formalization inventory. All essential
mathematical arguments used by the package are stated and proved here.

\section{Extensions and boundaries of the method}
\subsection{Other endpoints and nonunit leading powers}
A problem at infinity may sometimes be moved to zero by reciprocal
coordinates, but the transformed equation must still be checked for positive
excess exponents and a nonzero linear leading term. The direction of the
asymptotic ordering changes under such substitutions; it should not be
handled by merely reversing a sorting routine.

For a leading term \(a x^\beta\), \(a,\beta>0\), one can first set
\(t=(y/a)^{1/\beta}\). Taking the positive \(\beta\)-th root of the normalized
forward equation produces a near-identity equation. The root usually creates
an infinite binomial expansion rather than a finite exact log-power sum.
Finite determination and a controlled remainder, as in
Proposition~\ref{prop:stability}, then justify finite-jet inversion. That
preprocessing is a mathematical extension, not an automatic feature of this
release.

\subsection{Faster engines without changing the specification}
The multi-index coefficient formula favors transparency. At high cutoffs,
one could replace repeated polynomial products by cached powers, use a
priority queue for support enumeration, or use truncated Newton iteration on
\eqref{eq:correction}. Its derivative with respect to \(u\) is a unit
\(1+\) positive-valuation terms, so formal inversion is possible and Newton
steps increase the accuracy rapidly. These optimizations should preserve the
same exact cutoff and remainder semantics, and should still be checked by a
separate composition engine.

For many mutually dependent generators, a semigroup normal form or lattice
relation basis could avoid repeated representations. It must not discard
multiplicity: different representations contribute distinct terms in
\eqref{eq:master}. Any optimized enumeration must reproduce their sum.

\subsection{What fails outside the hypotheses}
If \(\alpha_j=0\), a logarithmic correction need not be small relative to the
linear term, and the positive-valuation iteration can fail. If infinitely many
positive generators accumulate at zero, local finite enumeration fails.
Negative logarithm powers require rational-function rather than polynomial
coefficients and different zero-weight bookkeeping. Oscillatory terms can
require an enlarged coefficient algebra and separate branch estimates.
These cases need genuinely additional hypotheses; accepting them silently
would undermine both termination and the analytic error claims.

\section{Conclusion}
The requested inverses do not require a Taylor series at the singular
endpoint. They require a branch condition and the right asymptotic algebra.
For finite positive log-power perturbations of a linear germ, the answer is
fully constructive: normalize the linear coefficient, enumerate a finite
positive exponent simplex, apply the Euler-polynomial Lagrange formula,
merge coincident weights, and report a logarithm-aware remainder. Analytic
continuation in an auxiliary parameter validates the formal answer and
provides an explicit conditional error majorant. The two examples are
special cases of the same theorem, with Catalan leading logarithms in one
case and generalized binomial coefficients on an irrational grid in the
other.

\appendix
\section{Archive contents and provenance}
\begin{center}
\begin{tabular}{@{}p{0.46\linewidth}p{0.47\linewidth}@{}}
\toprule
Path & Purpose\\\midrule
\texttt{article/article.tex}, \texttt{article.pdf} & This source and compiled article.\\
\texttt{Kernel/RealLogPowerInverse.wl} & Main Wolfram Language package.\\
\texttt{Kernel/init.m}, \texttt{PacletInfo.wl} & Loading and package metadata.\\
\texttt{Examples/Examples.wl} & Executable usage examples.\\
\texttt{Tests/RealLogPowerInverse.wlt} & Native regression tests, not executed here.\\
\texttt{Tests/RunTests.wls} & Command-line native test runner.\\
\texttt{validation/reference\_validation.py} & Executed independent symbolic/numerical checks.\\
\texttt{validation/results.json}, \texttt{run.log} & Recorded independent results.\\
\texttt{README.md}, \texttt{VALIDATION.md} & Usage and explicit testing status.\\
\texttt{LICENSE}, \texttt{SHA256SUMS} & License for newly written code; checksums.\\
\bottomrule
\end{tabular}
\end{center}
The user's source archive was consulted but is not repackaged here. Public
sources are cited for attribution and context. The mathematical proofs,
package, and independent test implementation in this archive were prepared
for this task. The code license does not purport to relicense any cited
third-party material.

\begin{thebibliography}{9}
\bibitem{question}
V. Reshetnikov, \emph{How to get an asymptotic of the real-valued branch of the
inverse function?}, Mathematica Stack Exchange, question 236367,
December 13, 2020; with answers by Michael E2 and Somos.
\url{https://mathematica.stackexchange.com/questions/236367/}.
The user-supplied archive and the public page were consulted.

\bibitem{dlmf}
NIST Digital Library of Mathematical Functions,
\S1.10(vii), \emph{Inverse Functions}, especially the Lagrange inversion
formulas. \url{https://dlmf.nist.gov/1.10#vii}.

\bibitem{proveit}
V. Reshetnikov and contributors, \emph{ProveIt}, polynomial-logarithmic
series and transseries materials, commit
\texttt{24ce8bd743eaab64a91ce90725ea00f498d319d2}.
\url{https://github.com/VladimirReshetnikov/ProveIt/tree/24ce8bd743eaab64a91ce90725ea00f498d319d2/Analysis/FabiusFunction/docs/semi-formalized-research-frontiers/drafts/series-and-transseries}.
See the source-inspection limitation in Section~\ref{sec:related}.

\bibitem{inverseseries}
Wolfram Research, \emph{InverseSeries}, Wolfram Language documentation.
\url{https://reference.wolfram.com/language/ref/InverseSeries.html}.

\bibitem{rootreduce}
Wolfram Research, \emph{RootReduce}, Wolfram Language documentation.
\url{https://reference.wolfram.com/language/ref/RootReduce.html}.

\bibitem{algebraics}
Wolfram Research, \emph{Algebraics}, Wolfram Language documentation.
\url{https://reference.wolfram.com/language/ref/Algebraics.html}.

\bibitem{testreport}
Wolfram Research, \emph{TestReport}, Wolfram Language documentation.
\url{https://reference.wolfram.com/language/ref/TestReport.html}.
\end{thebibliography}
\noindent\small Online sources were consulted on September 7, 2026.
\end{document}
